\section{Angular continuation of the compact zeta difference}\label{app:zeta}
Let $p_k=m_+(k)-m_-(k)$ denote shell differences. A finite-group character average gives the multiplicities \cite{BarSpin}. For a left binary group,
\[
 m_+(k)=(k+2)\dim(V_k\otimes\chi)^\Gamma,\quad
 m_-(k)=(k+1)\dim(V_{k+1}\otimes\chi)^\Gamma .
\]
The identity contribution cancels in $p_k$. The central $-1$ contributes an alternating quadratic. The other conjugacy classes contribute a bounded periodic character $\sin((k+1)\alpha)/\sin\alpha$ times a factor linear in $k$. For the cyclic two-sided cases one counts pairs of allowed weights or uses the corresponding two-sided character products. Thus $p_k=\sum_{j=0}^2 c_j(k)k^j$ with periodic coefficients. For every line of Table~\ref{tab:eta} the three averages $\sum_{r=0}^{P_0-1}c_j(r)$ vanish, as verified with exact rational arithmetic for the indicated period $P_0$.

A practical formula follows by setting $\sum p_k x^k=N(x)/(1-x^{P_0})^3$. Its numerator is obtained by three differences in steps of $P_0$ and has degree less than $3P_0$. Three integer moment sums give $N(1)=N'(1)=N''(1)=0$. The limit is $-N'''(1)/(6P_0^3)$. A Mellin transform of $e^{-3t/2}\sum p_k e^{-kt}$ gives the same value for the angular eta function at zero. For instance the even $\mathbb{RP}^3$ shell sequence is $(-1)^k(k+1)(k+2)$, with generating function $2/(1+x)^3$ and limit $1/4$. This evaluates a global angular invariant without truncating an oscillatory sum.

\paragraph{Agreement with four-dimensional continuation.}
Let $Z_K(z)$ be the difference of the two positive-shell zeta functions of $A_s^*A_s$, $K=k+3/2$. Opposite signs of $\lambda$ exchange adjoint and charge, so the direct sum is initially $\sum_k p_k Z_K(z)$ for $\Re z>2$. At each fixed shell $Z_K(0)=0$ and $-Z'_K(0)=-2J$. We justify taking the finite part of this divergent angular sum.

Extend large $K$ continuously. The first-order self-adjoint radial Dirac operator is affine in $K$ with bounded slope on the fixed compact collar and $K$-independent separated domain. For either square, the potential is $b_s^2\mp b_s'$, bounded below by $cK^2$ for large $K$; its Robin form has a nonnegative endpoint contribution in the corresponding block. The min--max principle and the mixed-interval Laplacian yield a lower bound $c'(K^2+n^2)$ for squared radial eigenvalues, uniformly in $K$ (adjusting constants for finitely many levels). A resolvent contour outside the two real spectral rays can use fixed angles and inner radius proportional to $K$. There the operator norm is $O((K+|\xi|)^{-1})$ and the square of the Hilbert--Schmidt norm is $O((K+|\xi|)^{-1})$, by summing the last eigenvalue estimate. After $m$ differentiations in $K$, the resolvent product has $m+1$ factors and bounded insertions. Two factors in the Hilbert--Schmidt norm and the rest in operator norm bound its trace by $O((K+|\xi|)^{-m})$. Integrating $(\pm\xi)^{-2z}$ on the spectral contours therefore gives locally uniformly
\[
 |\partial_K^m Z_K(z)|\le C_z K^{1-m-2\Re z},\qquad m=4,5,6,\quad\Re z>-1/2.
\]
The contour can be held fixed in a neighborhood of each $K$ while differentiating. Only these differentiated traces are estimated close to $z=0$.

Each zero-mean periodic sequence admits a bounded zero-mean discrete primitive, iteratively six times. Sum by parts six times for $c_j(k)k^j Z_K(z)$. The sixth difference of $k^j Z_K$ is $O(K^{1+j-6-2\Re z})$, using the three derivative bounds above. This is summable at the origin and locally uniformly near it, and the finite initial boundary terms survive. Differentiation there leaves sixth differences of a degree-at-most-two polynomial times the constant $-2J$, hence a zero remainder. To evaluate the initial terms, perform the same summation by parts on $k^j\rho^k$ and take $\rho\uparrow1$. The absolute sum of sixth differences tends to zero as $O((1-\rho)^{5-j})$. The surviving coefficient is the angular Abel sum computed above, proving \eqref{eq:delta}. The zero of $Z_K(0)$ likewise removes a common determinant scale from the ratio. This cancellation relies on the structure of the \emph{difference} of signed branch spectra.
